\documentclass[10pt,a4paper]{amsart}
\usepackage[margin=19mm]{geometry}
\usepackage{amsmath,amssymb,mathtools}
\usepackage[scr=rsfs,cal=euler]{mathalfa}
\usepackage{xfrac}
\usepackage[unicode,hidelinks,bookmarks=false]{hyperref}
\newcommand{\ie}{i.e.\ }
\newcommand{\cf}{cf.\ }
\newcommand{\resp}{respectively\ }
\newcommand{\Subsection}{\subsection}
\providecommand{\nobreakdash}{\nobreak\hbox{-}}
\setlength{\emergencystretch}{2em}
\multlinegap0pt
\usepackage{bm}
\usepackage{bbm}
\usepackage{dsfont}
\usepackage{xspace}
\usepackage{cancel}
\usepackage{mathtools}
\usepackage{amsthm}
\makeatletter
\@ifundefined{lemma}{\newtheorem{lemma}{Lemma}}{}
\@ifundefined{proposition}{\newtheorem{proposition}[lemma]{Proposition}}{}
\@ifundefined{theorem}{\newtheorem{theorem}[lemma]{Theorem}}{}
\makeatother
\newcommand{\e}{\varepsilon} 
\title[Non-local singular perturbations of non-convex functionals -]{Non-local singular perturbations of non-convex functionals -- recent results}
\author{Andrea Braides}
\date{Source: arXiv:2601.08573v1 (CC BY 4.0)}
\begin{document}
\maketitle

\noindent\emph{Source and licence.} Non-local singular perturbations of non-convex functionals -- recent results. Version arXiv:2601.08573v1 (CC BY 4.0), distributed under the Creative Commons Attribution 4.0 International licence, as recorded in the arXiv licence metadata for this exact version and shown on the abstract page https://arxiv.org/abs/2601.08573v1. Full rights notice: licenses/arxiv-2601.08573-CC-BY-4.0.txt.\par\smallskip

\noindent\textit{Editorial scope.} Counted unit: the Theorem whose source label is \texttt{k+s-lim-1} in the article (source0.tex lines 232--238, source label \texttt{k+s-lim-1}), delivered together with the article's own complete original proof of it (source0.tex lines 240--271, one \texttt{proof} environment of 32 lines). No mathematics is written, repaired or re-derived: statement and proof are the article's own text line for line. Required context retained uncounted: sut (lines 185--194). Every definition, display and label of those lines is retained unchanged. Omitted from this package and not redistributed: 1--184, 195--231, 239--239, 272--543. Nothing inside a retained excerpt is omitted or altered: every retained body is delivered exactly as the article's lines are written, so no omission receipt of any kind accompanies this package. Citations: the retained excerpts carry no citation command at all, so no bibliography entry of the article is transcribed and no reference is redistributed. Reference adaptation: none was needed and none was made; every reference inside the delivered unit resolves inside it, so no unresolved reference or citation is produced. Printed slips kept literal: no printed word, symbol, hypothesis or number of the counted statement or of its proof was altered. Layout and class: the article's own class is \texttt{article} with packages including amsmath, amsthm, amssymb, color, xcolor; the wrapper is the lane's pinned \texttt{amsart} editorial wrapper at 10pt on A4, which restates the theorem-like environments the retained text needs and adds the article's own macro definitions that the retained text needs (\texttt{\textbackslash e}), with extra wrapper packages (bm, bbm, dsfont, xspace, cancel, mathtools). The delivered numbering is the unit's own. Assets: no retained span contains a figure environment, a graphics-inclusion command, a PDF-page inclusion, a table or a TikZ picture; no image file is copied into this package, the package declares no asset dependency and no image file is redistributed with it. Licence: version arXiv:2601.08573v1 of this article is distributed under the Creative Commons Attribution 4.0 International licence, as recorded in the arXiv licence metadata for this exact version and shown on the abstract page https://arxiv.org/abs/2601.08573v1. A full rights notice is delivered at licenses/arxiv-2601.08573-CC-BY-4.0.txt. Characters: every retained excerpt is pure ASCII, so the delivered engine, XeLaTeX with its default Latin Modern text font, reports no missing character.
\begin{proposition}[$k+s$-surface tension]\label{sut}
Let $k\in \mathbb N$ and $s\in (0,1)$ with $k+s>\frac12$. Then we have the equality
\begin{eqnarray}\label{emme}\nonumber
&&m_{k+s} =\inf\Big\{\int_{\mathbb R} W(v)dx+\int_{\mathbb R^2}\frac{|v^{(k)}(x)-v^{(k)}(y)|^2}{|x-y|^{1+2s}} dx\,dy :\\
&& \hskip4.5cm v\in H^{k+s}_{\rm loc}(\mathbb R), \lim_{x\to\pm\infty}v(x)=\pm1\Big\}\\
\label{emmetilde}\nonumber
&&\hskip1cm =\inf_{T>0} \inf\bigg\{\int_{-T}^{+T} W(v)dt +\int_{\mathbb R^2}\frac{|v^{(k)}(x)-v^{(k)}(y)|^2}{|x-y|^{1+2s}} dx\,dy :\\
&& \hskip4cm v\in H^{k+s}_{\rm loc}(\mathbb R), v(x)={\rm sign}\, x\hbox{ \rm if }|x| \ge T\Big\}.
\end{eqnarray}
\end{proposition}

\begin{theorem}[sharp-interface limit] Let $k\in \mathbb N$ and $s\in (0,1)$ with $k+s>\frac12$; then
we have 
 \begin{equation}\label{k+s-lim-1}
F^{k+s}(u)= m_{k+s} \#(S(u)).
\end{equation}
for all $u\in BV(I;\{-1,1\})$.
\end{theorem}

\begin{proof} Let $u_\e\to u$ in measure with $F^{k+s}_\e(u_\e)\le S<+\infty$. Let $t_0\in S(u)$, and define 
$\widetilde u_\e(t) = u_\e(t_0+\e t)$. We have $F^{k+s}(\widetilde u_\e, (-T,T))\le S$ for all $T>0$, and argue similarly to the proof of  Proposition \rm\ref{sut}. Given $\eta>0$, we fix $x_0$ such that, up to a change of sign, $|\widetilde u_\e(x)-{\rm sign}(x)|\le \eta$ for $|x|\in[x_0, T]$. We can construct $v_\e\in H^{k+s}_{\rm loc}(\mathbb R)$ such that $v_\e(x)={\rm sign}(x)$ if $|x|\ge \frac23 T$ and for all $\delta>0$ there exists $C_\delta$ such that 
\begin{eqnarray*}\nonumber
F^{k+s}(v_\e,\mathbb R))\le (1+\delta) F^{k+s}_\e(u_\e,(t_0-\e T,t_0+\e T)) +\frac{C_\delta}{T^{2s}}.
\end{eqnarray*}
Since $v_\e$ is a test function for $m_{k+s}$, summing up for $t_0\in S(u)$ and taking into account that all $(t_0-\e T,t_0+\e T)$ are disjoint for $\e$ small enough, we obtain that
$$
m_{k+s}\,\#(S(u))\le  (1+\delta) \liminf_{\e\to 0} F^{k+s}_\e(u_\e) + \frac{C_\delta}{T^{2s}}.
$$
By the arbitrariness of $T$ and $\delta$ we obtain the liminf inequality. 

\smallskip
The construction of an approximate recovery sequence for a function $u\in BV(I;\{-1,1\})$ follows a usual pattern:
given $\eta>0$ we find $T>0$ and $v\in H^{k+s}_{\rm loc}(\mathbb R)$ such that $v(x)={\rm sign}\, x$ if $|x| \ge T$ and 
$$
F^{k+s}(v,\mathbb R)\le m_{k+s}+\eta.
$$
Let $S(u)=\{t_1,\ldots, t_N\}$ and define 
$$
u_\e(t)=\begin{cases}
v\Big(\frac{u(t_j^+)-u(t_j^-)}2\cdot\frac{t-t_j}\e\Big) & \hbox{ if } |t-t_j|\le \e\\
u(t) & \hbox{ otherwise.}
\end{cases}
$$
We set $x_0=\inf I$, $x_j=(t_{j+1}+t_j)/2$ for $j\in\{1,\ldots, N-1\}$ and $x_N=\sup I$.
In order to check that $u_\e$ is a recovery sequence, we need to show that the amount of energy due to
the interactions between $x$ and $y$ belonging to intervals $(x_{j_1-1}, x_{j_1})$ and $(x_{j_2-1}, x_{j_2})$ with $j_1\neq j_2$. If $k\ge 1$ we note that actually the interaction is only between the intervals $(t_{j_1}-\e T,t_{j_1}+\e T)$ and $(t_{j_2}-\e T,t_{j_2}+\e T)$, so that $|x-y| \sim |t_{j_2}-t_{j_1}|$ independently of $\e$, which gives that these contributions are asymptotically negligible. In the case $k=1$ one additionally has to note that for ${j_2}={j_1}+1$ there is no interaction between pair of points in $(t_{j_1}+\e T, t_{j_2}-\e T)$. This implies that
$$
F_\e^{k+s}(u_\e,I)\le (m_{k+s}+\eta)\#(S(u)) +o(1)
$$
as $\e\to 0$, which proves the claim by the arbitrariness of $\eta$.
\end{proof}

\end{document}
